\documentclass[11pt]{amsart}

\usepackage[T1]{fontenc}
\usepackage{lmodern}
\usepackage{microtype}
\usepackage{amsmath,amssymb,amsthm}
\usepackage[hidelinks]{hyperref}

\newtheorem{theorem}{Theorem}
\newtheorem{corollary}[theorem]{Corollary}

\hypersetup{
  pdftitle={A Counterexample to a Local Tor-Length Conjecture of Boix and Lima-Pereira}
}

\title[Local Tor-length counterexample]
{A Counterexample to a Local Tor-Length Conjecture of Boix and Lima--Pereira}

\date{}

\subjclass[2020]{Primary 13D02}
\keywords{Tor modules, Koszul homology, complete intersections, length inequalities}

\begin{document}

\begin{abstract}
We exhibit complete intersection ideals $I,J$ in a regular local ring of
dimension four such that
\[
\begin{gathered}
\ell(R/(I+J))=8,\\
\bigl(\ell(\operatorname{Tor}_i^R(R/I,R/J))\bigr)_{i=0}^3
=(8,22,22,8).
\end{gathered}
\]
The example disproves Conjecture~4.2 of Boix and Lima--Pereira and gives
negative answers to their Questions~5.1 and~5.4.
\end{abstract}

\maketitle

\section{The counterexample}

Write $\ell$ for module length. Boix and Lima--Pereira conjectured in
\cite[Conjecture~4.2]{BLP} that if
$(R,\mathfrak m)$ is a regular local ring of dimension $n$, $I$ is
generated by a regular sequence of length $k\leq n-1$, and $J$ is
generated by a regular sequence of length $n$, then
\begin{equation}\label{eq:conjecture}
\ell\bigl(\operatorname{Tor}_i^R(R/I,R/J)\bigr)
=
\binom{k}{i}\ell\bigl(R/(I+J)\bigr)
\qquad (0\leq i\leq k).
\end{equation}
They also asked whether the corresponding degreewise lower bound holds
when $I$ and $J$ merely contain regular sequences of lengths $k$ and $n$
\cite[Question~5.1]{BLP}, and whether
\begin{equation}\label{eq:total}
\sum_{i=0}^{n}\ell\bigl(\operatorname{Tor}_i^R(R/I,R/J)\bigr)
\geq 2^k\ell\bigl(R/(I+J)\bigr)
\end{equation}
holds under the same hypotheses \cite[Question~5.4]{BLP}.

\begin{theorem}\label{thm:main}
Let
\[
S=\mathbb Q[x,y,z,w],
\qquad
\mathfrak m=(x,y,z,w),
\qquad
R=S_{\mathfrak m}.
\]
Define ideals of $S$ by
\begin{align*}
I_S&=(x^2+xw,\ y^2+yw,\ z^2+zw),\\
J_S&=(x^2+xy,\ xz+z^2,\ y^2+yz,\ xw+w^2),
\end{align*}
and set $I=I_SR$ and $J=J_SR$. Then the displayed generators of $I$ and
$J$ are regular sequences of lengths $3$ and $4$, respectively, and
\[
\ell\bigl(R/(I+J)\bigr)=8.
\]
Moreover,
\[
\bigl(\ell(\operatorname{Tor}_i^R(R/I,R/J))\bigr)_{i=0}^3
=(8,22,22,8),
\]
and $\operatorname{Tor}_i^R(R/I,R/J)=0$ for $i\geq4$.
\end{theorem}

\begin{proof}
Successively quotienting by the three monic polynomials
$x^2+xw$, $y^2+yw$, and $z^2+zw$ shows that they form an $S$-regular
sequence. Moreover, $S/I_S$ is free over $\mathbb Q[w]$ with basis
\[
\{x^ay^bz^c:a,b,c\in\{0,1\}\}.
\]
Hence their images form an $R$-regular sequence.

We next show that $J_S$ is $\mathfrak m$-primary. Let $P$ be a prime
ideal containing $J_S$. If $x\in P$, then the second, third, and fourth
generators imply $z,y,w\in P$, respectively, so $P=\mathfrak m$.
If $x\notin P$, then $x+y\in P$. Thus $y\notin P$, and $y+z\in P$.
It follows that $z-x\in P$ and $z\notin P$. The relation
$z(x+z)\in P$ then gives $x+z\in P$, whence $2x\in P$, a contradiction.
Therefore $\sqrt{J_S}=\mathfrak m$. The four generators of $J$ form a
system of parameters in the Cohen--Macaulay local ring $R$, and hence an
$R$-regular sequence.

With graded reverse lexicographic order $x>y>z>w$, a Gr\"obner basis of
$I_S+J_S$ is
\begin{align*}
&x^2-w^2,\ xy+w^2,\ y^2+yw,\ xz-zw,\ yz-yw,\\
&z^2+zw,\ xw+w^2,\ w^2y,\ w^2z,\ w^3.
\end{align*}
Its standard monomials are
\[
1,\ x,\ y,\ z,\ w,\ w^2,\ yw,\ zw.
\]
Consequently, $I_S+J_S$ is $\mathfrak m$-primary and
\[
\ell_S\bigl(S/(I_S+J_S)\bigr)=8.
\]

A second exact Gr\"obner-basis computation gives
\begin{equation}\label{eq:hilbert-tor}
\operatorname{HS}_{(I_S\cap J_S)/(I_SJ_S)}(t)=8t^3+14t^4.
\end{equation}
The computation is reproduced below. Since
\[
\operatorname{Tor}_1^S(S/I_S,S/J_S)
\cong (I_S\cap J_S)/(I_SJ_S),
\]
we obtain
\[
\ell_S\bigl(\operatorname{Tor}_1^S(S/I_S,S/J_S)\bigr)=22.
\]
All modules in question are supported at $\mathfrak m$, so localization
at $\mathfrak m$ preserves their lengths and commutes with Tor. Thus
\[
\ell_R\bigl(\operatorname{Tor}_0^R(R/I,R/J)\bigr)=8,
\qquad
\ell_R\bigl(\operatorname{Tor}_1^R(R/I,R/J)\bigr)=22.
\]
Koszul duality gives
\[
\ell_R\bigl(\operatorname{Tor}_i^R(R/I,R/J)\bigr)
=
\ell_R\bigl(\operatorname{Tor}_{3-i}^R(R/I,R/J)\bigr)
\qquad (0\leq i\leq3);
\]
see \cite[Proposition~1.9]{BLP}. Hence the Tor-length vector is
$(8,22,22,8)$. Finally, $R/I$ has projective dimension $3$, so the higher
Tor modules vanish.
\end{proof}

\begin{corollary}\label{cor:consequences}
Conjecture~4.2 of \cite{BLP} is false. Questions~5.1 and~5.4 of
\cite{BLP} also have negative answers, even when both ideals are complete
intersections.
\end{corollary}

\begin{proof}
Here $n=4$, $k=3$, and $\ell(R/(I+J))=8$. For $i=1,2$,
\[
22<24=\binom{3}{i}\,8,
\]
which disproves \eqref{eq:conjecture} and the lower bound in
Question~5.1. Moreover,
\[
\sum_{i=0}^{4}\ell\bigl(\operatorname{Tor}_i^R(R/I,R/J)\bigr)
=8+22+22+8=60<64=2^3\cdot8,
\]
so \eqref{eq:total} fails.
\end{proof}

\section{Exact computation}

The following \textsc{Macaulay2} script \cite{M2} verifies
\eqref{eq:hilbert-tor} and the complete Tor-length vector over
$\mathbb Q$. Since all computed modules have finite length, the command
\texttt{degree} returns their vector-space dimensions.

{\footnotesize
\begin{verbatim}
S = QQ[x,y,z,w, MonomialOrder => GRevLex];

IS = ideal(x^2+x*w, y^2+y*w, z^2+z*w);
JS = ideal(x^2+x*y, x*z+z^2, y^2+y*z, x*w+w^2);

assert(codim IS == 3);
assert(codim JS == 4);
assert(apply(0..4, d -> hilbertFunction(d,S/(IS+JS)))
       == {1,4,3,0,0});

Q = (intersect(IS,JS))/(IS*JS);
assert(apply(0..5, d -> hilbertFunction(d,Q))
       == {0,0,0,8,14,0});
assert(degree Q == 22);

M = coker gens IS;
N = coker gens JS;
assert(apply(0..4, i -> degree(Tor_i(M,N)))
       == {8,22,22,8,0});
\end{verbatim}
}

\enlargethispage{4\baselineskip}

\begin{thebibliography}{9}

\bibitem{BLP}
A.~F. Boix and B.~K. Lima--Pereira,
\emph{Exploring a local variant of the
Buchsbaum--Eisenbud--Horrocks conjecture},
Proc. Roy. Soc. Edinburgh Sect. A, First View (2026), 1--20;
\href{https://doi.org/10.1017/prm.2026.10176}
{doi:10.1017/prm.2026.10176};
\href{https://arxiv.org/abs/2607.26118}{arXiv:2607.26118}.

\bibitem{M2}
D.~R. Grayson and M.~E. Stillman,
\emph{Macaulay2, a software system for research in algebraic geometry},
available at \url{https://macaulay2.com/}.

\end{thebibliography}

\end{document}